\clearpage                                  % В том числе гарантирует, что список литературы в оглавлении будет с правильным номером страницы
%\hypersetup{ urlcolor=black }               % Ссылки делаем чёрными
%\providecommand*{\BibDash}{}                % В стилях ugost2008 отключаем использование тире как разделителя
\urlstyle{rm}                               % ссылки URL обычным шрифтом
\ifdefmacro{\microtypesetup}{\microtypesetup{protrusion=false}}{} % не рекомендуется применять пакет микротипографики к автоматически генерируемому списку литературы
% \insertbibliofull                           % Подключаем Bib-базы: все статьи единым списком
% Режим с подсписками
\insertbiblioexternal                      % Подключаем Bib-базы: статьи, не являющиеся статьями автора по теме диссертации
% Для вывода выберите и расскомментируйте одно из двух
%\insertbiblioauthor                        % Подключаем Bib-базы: работы автора единым списком 
\insertbiblioauthorgrouped                 % Подключаем Bib-базы: работы автора сгруппированные (ВАК, WoS, Scopus и т.д.)
% Счётный проход окружения `counter' (biblio/biblatex.tex) наполняет счётчик citenum
% для фразы <<Список литературы содержит N наименований>> во введении.
% Идти он должен строго ПОСЛЕ всех видимых \printbibliography: при defernumbers=true
% номер присваивается записи при первом выводе, и более ранний проход по keyword=bibliofull
% пронумеровал бы работы автора в порядке .bbl, а не в порядке групп.
% Вывод записей подавлен драйвером окружения, но вертикальный материал оно всё равно
% порождает, поэтому собираем проход в отбрасываемый бокс --- иначе в печать может уйти
% лишняя пустая страница.
\ifnumequal{\value{draft}}{0}{\setbox0=\vbox{\printbibliography[heading=nobibheading,env=counter,keyword=bibliofull,section=0]}}{}
\ifdefmacro{\microtypesetup}{\microtypesetup{protrusion=true}}{}
\urlstyle{tt}                               % возвращаем установки шрифта ссылок URL
%\hypersetup{ urlcolor={urlcolor} }          % Восстанавливаем цвет ссылок
